\chapter{IMPLEMENTATION DETAILS}
\section{Hardware Implementation}
\subsection{pH Sensing (Commercial Pen-Type pH Meter)}
The pH measurement subsystem employs a commercial pen-type pH meter, whose amplified analog voltage output is connected to ESP32 GPIO34. GPIO34 is an input only pin belonging to the ADC1 group, which remains stable and unaffected during active Wi-Fi operation, making it well-suited for reliable analog sensing in this system.\\[8pt]
Rather than relying on the raw electrode voltage directly, the system performs a three-point calibration by immersing the pen-type pH meter into standard buffer solutions of pH 4.00, 6.86, and 9.18. The corresponding amplified voltage outputs from the pH module are recorded at each calibration point. These three voltage-to-pH pairs establish a piecewise linear mapping across two segments, one from pH 4.00 to 6.86 and another from pH 6.86 to 9.18, allowing the firmware to interpolate and convert any measured voltage into its corresponding pH value with improved accuracy across the operational range of the system. Linearly interpolating between the datapoints to calculate constants forming a line of best fit describing the relationship between the electrode voltage and the actual pH of the solutions, the mid-pH i.e. measured electrode voltage at pH 7, is used as such in the following equation to convert voltage readings from the ESP32s ADC to corresponding pH:
\begin{equation}
pH = 7 - \frac{MID\_PH_v - pH\_raw_v}{CALIBRATION\_FACTOR}
\end{equation}
where,\\
$pH$ = calculated acidity or alkalinity of the solution, \\
$MID\_PH_v$ = reference midpoint value in volts at ph 7, \\
$pH\_raw_v$ = raw sensor reading in volts,\\
$CALIBRATION\_FACTOR$ = constant used for calibration.

The pen-type meter provides a stable, pre-buffered voltage signal to the module, and the ESP32 samples this through its 12-bit ADC configured with full 3.3V attenuation, yielding a resolution of approximately 0.01 pH units.

\subsection{TDS Meter (TDS Meter V1.0 Module)}
The TDS measurement system connects the conductivity probe to the TDS Meter V1.0 module, which outputs a 0-3.3V analog signal directly to ESP32 GPIO35. GPIO35 is an input-only pin from the ADC1 bank, ensuring stable analog readings even while Wi-Fi is actively operating. The TDS module internally generates an approximately 1 kHz square wave signal to prevent electrode polarization, then converts the resulting AC signal to a DC voltage proportional to the solution's conductivity. The sampling protocol incorporates multiple rapid readings with averaging to mitigate electrical noise inherent in conductivity measurements.\\[8pt]
In the integrated monitoring system, the TDS sensor measures Total Dissolved Solids concentration in parts per million, employing electrical conductivity measurement converted to TDS through established empirical relationships. The sensor exhibits significant temperature dependence with a 2\% per degree Celsius coefficient, making DHT22 temperature readings essential for accurate compensation.

\subsection{DHT22 Temperature and Humidity Sensor}
The DHT22 sensor interfaces with the ESP32 through a single-wire digital protocol connected to GPIO32, with an external 4.7 k$\Omega$ pull-up resistor to 3.3V. GPIO32 is a fully capable general-purpose digital I/O pin from the ADC1 group and is free from boot-strapping constraints or RF subsystem conflicts, ensuring reliable single-wire communication. The sensor operates from a 3.3V supply with recommended decoupling capacitors placed near the sensor to ensure stable operation. The DHT22 arrives from the factory with pre-loaded calibration coefficients, typically eliminating the need for user calibration in standard applications.\\[8pt]
Within the integrated sensor system, the DHT22 serves dual critical functions. Primarily, it provides accurate ambient temperature and relative humidity measurements with specified accuracies of ±0.5°C across a -40°C to 80°C range and ±2-5\% across 0-100\% relative humidity. Secondarily and equally importantly, it supplies the essential temperature data required for compensating both pH and TDS measurements. This cross-sensor dependency establishes the DHT22 as a foundational component whose readings directly impact the accuracy of the other two sensors.

\section{Software Implementation}
The software follows a client-server model which makes the system flexible. It consists of firmware running on the ESP32, a backend server, a cloud database, and a frontend client application.

\subsection{ESP32 Firmware}
The ESP32 firmware is responsible for sensor data acquisition, basic preprocessing, Wi-Fi communication and transmission of data. The firmware is such that it is lightweight and efficient so that the firmware remains stable even in the face of constant monitoring. The sensor data is occasionally retrieved and sent to the cloud database via the HTTP-based communication.

\subsection{Backend Server Architecture}
The backend of the system is created under the Django web framework and Django Rest Framework (DRF). The backend serves as an intermediary between the database and the frontend application. It offers RESTful API endpoints to data ingestion, data retrieval, system authentication and system configuration.\\[8pt]
To guarantee a secure system access, JWT-based authentication is carried out. There is management of Cross-Origin Resource Sharing (CORS) that facilitates a smooth flow between the frontend and the backend.

\subsection{Database Architecture}
Supabase has been employed to store sensor data with the help of a cloud-based database. It offers a scalable and dependable data storage and retrieval platform which is real time. Sensors (pH, TDS, temperature, and humidity) data is stored with the timestamps to analyze the past and evaluate the performance.\\[8pt]
The cloud database will make this data persistent, accessible across the network, and compatible with the optimization methods of the future that will be based on data.

\subsection{Frontend Architecture}
The frontend application is developed using React and React Native to provide responsive and intuitive user interface. The frontend communicates with the backend server through RESTful APIs using Axios. It displays real time sensor readings, historical trends and system status through a dashboard-based layout. 


\section{Nutrient Solution}
To prepare the nutrient solution for the project, the necessary materials were gathered, including 120 grams of NPK 19-19-19, 180 grams of Calcium Nitrate, 60 grams of Epsom Salt, and two 500ml bottles of distilled water. Warm to hot distilled water was utilized to facilitate the rapid and complete dissolution of the solid components.
The preparation was conducted in two separate stages to prevent nutrient precipitation. For the creation of Solution A, 120 grams of NPK were dissolved into 250ml of warm distilled water, while 60 grams of Epsom Salt were dissolved into a separate 250ml volume of warm distilled water. Once it was confirmed that all particles were completely dissolved, these two liquids were combined into a single concentrate.\\[8pt]
Solution B was prepared by dissolving 180 grams of Calcium Nitrate into the remaining 500ml of warm distilled water. The mixture was stirred until the solution achieved full clarity, ensuring the concentrate was stable and ready for dilution.
.